\documentclass{article}
\usepackage[preprint]{staix_2026}

\usepackage[utf8]{inputenc}
\usepackage[T1]{fontenc}
\usepackage{hyperref}
\usepackage{url}
\usepackage{booktabs}
\usepackage{amsfonts}
\usepackage{nicefrac}
\usepackage{microtype}
\usepackage{xcolor}
\usepackage{graphicx}

\title{Formatting Instructions for STAI-X 2026 Submissions}

\author{STAI-X 2026 Organizing Committee}

\begin{document}

\maketitle

\begin{abstract}
This document describes the official formatting requirements for submissions to the First International Conference on Statistics and Trustworthy AI for Cross (X)-Domain Acceleration
 (STAI-X 2026). The conference has two tracks: a Poster Track with a 2-page main-content limit, and a Full Paper Track with an 8-page main-content limit, both including figures and tables. References do not count toward the main-content limit.
\end{abstract}

\section{Submission of papers to STAI-X 2026}
STAI-X 2026 accepts electronic submissions through the conference submission system. Authors should prepare their manuscripts using the official \LaTeX{} style files. Please use the current version of the template and do not modify the style file except where explicitly allowed.

\subsection{Tracks}
STAI-X 2026 has two submission tracks:
\begin{itemize}
    \item \textbf{Poster Track.} The main content is limited to \textbf{2 pages} including figures and tables, excluding references.
    \item \textbf{Full Paper Track.} The main content is limited to \textbf{8 pages} including figures and tables, excluding references.
\end{itemize}
Optional appendices may follow the references unless the call for papers states otherwise.

\subsection{Style}
Papers submitted to STAI-X 2026 must be prepared according to the instructions in this document. The official style file supports anonymous submissions, final camera-ready versions, and non-anonymous preprints.

\subsection{Retrieval of style files}
The official STAI-X 2026 template package consists of the following files:
\begin{itemize}
    \item \verb|staix_2026.sty| --- the conference style file;
    \item \verb|staix_2026.tex| --- an author shell for submissions;
    \item \verb|staix_2026_instructions.tex| --- this instruction document.
\end{itemize}

The style file supports the following options:
\begin{itemize}
    \item \verb|poster| --- Poster Track;
    \item \verb|full| --- Full Paper Track (default);
    \item \verb|final| --- camera-ready mode;
    \item \verb|preprint| --- non-anonymous preprint mode;
    \item \verb|nonatbib| --- disables automatic loading of \verb|natbib|.
\end{itemize}

Typical usage examples are:
\begin{verbatim}
\usepackage{staix_2026}
\usepackage[poster]{staix_2026}
\usepackage[final]{staix_2026}
\usepackage[poster,final]{staix_2026}
\usepackage[preprint]{staix_2026}
\end{verbatim}

\section{General formatting instructions}
\label{gen_inst}
The text must be confined within a rectangle 5.5 inches wide and 9 inches long. The left margin is 1.5 inches. Use 10-point type with 11-point leading. The title is centered between two horizontal rules. Paragraphs are separated by approximately half a line space, with no indentation.

In anonymous submission mode, author names are suppressed and line numbers are added automatically to aid review. In final and preprint modes, author information is shown and line numbers are removed.

\section{Headings: first level}
\label{headings}
All headings should be flush left and bold. First-level headings should be in 12-point type.

\subsection{Headings: second level}
Second-level headings should be in 10-point type.

\subsubsection{Headings: third level}
Third-level headings should be in 10-point type.

\section{Citations, figures, tables, references}
\label{others}

\subsection{Citations within the text}
Citations should be handled consistently. If \verb|natbib| is used, parenthetical citations can be created with \verb|\citep| and textual citations with \verb|\citet|.

\subsection{Footnotes}
Footnotes should be used sparingly. Place them at the bottom of the page on which they appear.

\subsection{Figures}
Figures should be legible and accompanied by captions placed below the figure.

\begin{figure}[h]
\centering
\fbox{\rule[-.5cm]{0cm}{4cm} \rule[-.5cm]{4cm}{0cm}}
\caption{Sample figure caption.}
\end{figure}

\subsection{Tables}
Tables should be legible and accompanied by captions placed above the table.

\begin{table}[h]
\caption{Sample table title}
\centering
\begin{tabular}{ll}
\toprule
Part & Description \\
\midrule
Style file & Controls formatting \\
Author shell & Starting point for submissions \\
Instruction file & Explains the rules \\
\bottomrule
\end{tabular}
\end{table}

\subsection{References}
References should appear after the main text. References do not count toward the main-content page limit for either track.

\section{Anonymity, preprints, and camera-ready versions}
At submission time, authors should omit the \verb|final| and \verb|preprint| options. This creates an anonymous submission with line numbers. Accepted papers should use the \verb|final| option for the camera-ready version. Authors who wish to post a non-anonymous preprint may use the \verb|preprint| option.

\section{Track-specific page limits}
The main-content page limits are:
\begin{itemize}
    \item \textbf{Poster Track:} 2 pages, excluding references.
    \item \textbf{Full Paper Track:} 8 pages, excluding references.
\end{itemize}
These limits apply to the main body of the paper, including figures and tables that appear within the main body.

\section{Acknowledgments and appendices}
Acknowledgments should not appear in the anonymous submission if they reveal author identity. In the supplied style file, the \verb|ack| environment is automatically hidden in anonymous submission mode and shown in final/preprint modes.

Appendices, proofs, and additional details may follow the references unless the call for papers specifies otherwise.

\bibliographystyle{plainnat}
\begin{thebibliography}{9}
\bibitem[Goodfellow et~al.(2016)Goodfellow, Bengio, and Courville]{goodfellow2016deep}
Ian Goodfellow, Yoshua Bengio, and Aaron Courville.
\newblock \emph{Deep Learning}.
\newblock MIT Press, 2016.

\bibitem[Lamport(1994)]{lamport1994latex}
Leslie Lamport.
\newblock \emph{\LaTeX: A Document Preparation System}.
\newblock Addison-Wesley, 2nd edition, 1994.
\end{thebibliography}

\end{document}
